\documentclass[aspectratio=169]{beamer}
\usetheme{Boadilla}
%Beamerフォント設定
\usefonttheme{professionalfonts}
\usepackage[T1]{fontenc}
\usepackage{mlmodern}  % 太いComputer Modern
% MLmodernのバグを修正: cf. https://tex.stackexchange.com/questions/646333/size-of-integral-symbol-in-section-header-with-mlmodern
\DeclareFontFamily{OMX}{mlmex}{}
\DeclareFontShape{OMX}{mlmex}{m}{n}{%
   <->mlmex10%
   }{}
\usepackage{newtxtext} % 数式以外をTXフォントで上書き
\renewcommand{\familydefault}{\sfdefault}  % 英文をサンセリフ体に
\usefonttheme{structurebold} % タイトル部を太字
\setbeamerfont{alerted text}{series=\bfseries} % Alertを太字
\setbeamerfont{section in toc}{series=\mdseries} % 目次は太字にしない
\setbeamerfont{frametitle}{size=\Large} % フレームタイトル文字サイズ
\setbeamerfont{title}{size=\LARGE} % タイトル文字サイズ
\setbeamerfont{date}{size=\small}  % 日付文字サイズ

%Beamer色設定
\definecolor{UniBlue}{RGB}{0,150,200}
\definecolor{UniGreen}{RGB}{3,175,122}
\definecolor{UniPink}{RGB}{255,128,255}
\definecolor{AlertOrange}{RGB}{255,76,0}
\definecolor{AlmostBlack}{RGB}{38,38,38}
\setbeamercolor{normal text}{fg=AlmostBlack}  % 本文カラー
\setbeamercolor{structure}{fg=UniBlue} % 見出しカラー
\setbeamercolor{block title}{fg=UniBlue!50!black} % ブロック部分タイトルカラー
\setbeamercolor{alerted text}{fg=AlertOrange} % \alert 文字カラー
\setbeamercolor{bibliography entry author}{fg=AlmostBlack} % Bib 文字カラー
\setbeamercolor{bibliography entry note}{fg=AlmostBlack} % Bib 文字カラー
\mode<beamer>{
    \definecolor{BackGroundGray}{RGB}{254,254,254}
    \setbeamercolor{background canvas}{bg=BackGroundGray} % スライドモードのみ背景をわずかにグレーにする
}

%フラットデザイン化
\setbeamertemplate{blocks}[default] % タイトルと本文を単色の背景でつなぐ
\useinnertheme{circles} % 箇条書きをシンプルに
\setbeamertemplate{navigation symbols}{} % ナビゲーションシンボルを消す
\setbeamertemplate{footline}[frame number] % フッターはスライド番号のみ

%タイトルページ
\setbeamertemplate{title page}{%
    \vspace{2.5em}
    {\usebeamerfont{title} \usebeamercolor[fg]{title} \inserttitle \par}
    {\usebeamerfont{subtitle}\usebeamercolor[fg]{subtitle}\insertsubtitle \par}

    \vspace{1.5em}
    \begin{columns}[]
        \begin{column}{.35\linewidth}
        \usebeamerfont{titlegraphic}\inserttitlegraphic\par
        \end{column}
        \begin{column}{.6\linewidth}\setlength\topsep{0pt}
            \begin{flushright}
        \usebeamerfont{author}\insertauthor\par
        \usebeamerfont{institute}\insertinstitute \par
        \vspace{3em}
        \usebeamerfont{date}\insertdate\par
            \end{flushright}
        \end{column}
    \end{columns}
}


% Algorithm系
\usepackage{algorithm}
\usepackage[noend]{algorithmic}
\algsetup{linenosize=\color{fg!50}\footnotesize}
\renewcommand\algorithmicdo{:}
\renewcommand\algorithmicthen{:}
\renewcommand\algorithmicelse{\textbf{else}:}
\renewcommand\algorithmicrequire{\textbf{Input:}}
\renewcommand\algorithmicensure{\textbf{Output:}}
\renewcommand\algorithmiccomment[1]{\hfill\textcolor{fg!50}{\footnotesize\textit{// #1}}}

% TikZ
\usepackage{tikz}
\usetikzlibrary{positioning,shapes,arrows,quotes,shapes.callouts,calc,graphs,graphs.standard,fit,backgrounds,perspective,matrix,decorations.pathmorphing,decorations.pathreplacing,decorations.shapes}
\tikzset{point/.style={circle, fill, inner sep=1pt, text width=1pt}}
\tikzset{alt/.code args={<#1>#2#3}{\alt<#1>{\pgfkeysalso{#2}}{\pgfkeysalso{#3}}}}
\tikzset{>=latex}
\tikzset{mynodes/.style={circle,white,fill=UniBlue!90,text width=.8em,inner sep=0pt,text centered,font=\footnotesize}}
\tikzset{myedges/.style={thick}}
\tikzset{mypath/.style={ultra thick, draw=#1}, mypath/.default=AlertOrange}
\tikzset{mysubset/.style={draw,#1,dotted,thick,fill=#1!10}, mysubset/.default=UniBlue}
\tikzset{mypath2/.style={ultra thick, dashed, draw=UniGreen}}
\tikzset{graphs/mygraph/.style={nodes={mynodes}, edges={myedges}, empty nodes}}
\tikzset{nomargin/.style={inner sep=0pt, outer sep=0pt}}
%% graph macros
\tikzgraphsset{declare={mybipartite}{%
subgraph I_nm [V={1,2,3,4,5,6}, W={a,b,c,d,e,f}];
1 -- {a,b};
2 -- {b,c,d,e};
3 -- {d,f};
4 -- {e,f};
{5,6} -- f;
}}
\tikzgraphsset{declare={mybipartitedi}{%
subgraph I_nm [V={1,2,3,4,5,6}, W={a,b,c,d,e,f}];
1 -> {a,b};
2 -> {b,c,d,e};
3 -> {d,f};
4 -> {e,f};
{5,6} -> f;
}}
\tikzgraphsset{declare={mybipartite2}{%
subgraph I_nm [V={1,2,3}, W={a,b,c}];
1 -- {a,b}; 2 -- {a,b,c}; 3 -- {a,c};
}} 
% Petersen graph
\tikzgraphsset{declare={Petersen}{%
    subgraph C_n [n=5,name=A, radius=1.5cm]; 
    subgraph I_n [n=5,name=B, radius=.8cm];
    \foreach \i [evaluate={\j=int(mod(\i+2,5)+1);}] in {1,...,5}{
        A \i -- B \i;
        B \i -- B \j;
}}}
% Example graph for Menger's theorem
\tikzgraphsset{declare={menger}{%
    s[at={(-1,0)}, as={$s$}];
    (s) -> 1[at={(0,1)}] -> 2[at={(1,0)}];
    (s) -> 3[at={(0,-1)}] -> 2;
    4[at={(2,1)}] <- 5[at={(2,-1)}];
    {5,6[at={(3,1)}], 7[at={(3,-1)}]} -> t[at={(4,0)}, as={$t$}];
    (s) -> (2);
    (2) -> {(4), (5)};
    {(4),(5)} -> {(6),(7)};
    (6) -> (7);
    1 -> 4;
}} 
\tikzgraphsset{declare={undirected menger}{%
    s[at={(-1,0)}, as={$s$}];
    (s) -- 1[at={(0,1)}] -- 2[at={(1,0)}];
    (s) -- 3[at={(0,-1)}] -- 2;
    4[at={(2,1)}] -- 5[at={(2,-1)}];
    {5,6[at={(3,1)}], 7[at={(3,-1)}]} -- t[at={(4,0)}, as={$t$}];
    (s) -- (2);
    (2) -- {(4), (5)};
    {(4),(5)} -- {(6),(7)};
    (6) -- (7);
    1 -- 4;
}} 

\usepackage{beamerweb}
\title{Beamer SVG Demo}
\author{Tasuku Soma}
\date{\today}

\begin{document}
\maketitle

\begin{frame}{HTML slider}
  \begin{block}{A Beamer block}
    \webid{formula}{$\displaystyle f(\lambda)=1-e^{-\lambda}$}
  \end{block}
  \webwidget[width=8cm,height=2cm,fallback={\texttt{Static value:} $\lambda=0.43$}]{slider-demo}
  \pause

  \only<2->{Move the slider to change the formula opacity.}
\end{frame}

\begin{frame}{TikZ network}
  \centering
  \begin{webinteractive}{flow-graph}
    \begin{tikzpicture}[x=1.7cm,y=1.3cm,thick]
      \begin{scope}[web id=vertex-s] \node[mynodes](s) at (0,0) {$s$}; \end{scope}
      \begin{scope}[web id=vertex-v] \node[mynodes](v) at (2,1) {$v$}; \end{scope}
      \begin{scope}[web id=vertex-w] \node[mynodes](w) at (2,-1) {$w$}; \end{scope}
      \begin{scope}[web id=vertex-t] \node[mynodes](t) at (4,0) {$t$}; \end{scope}
      \begin{scope}[web id=edge-sv] \draw[myedges,->] (s) -- (v); \end{scope}
      \begin{scope}[web id=edge-sw] \draw[myedges,->] (s) -- (w); \end{scope}
      \begin{scope}[web id=edge-vt] \draw[myedges,->] (v) -- (t); \end{scope}
      \begin{scope}[web id=edge-wt] \draw[myedges,->] (w) -- (t); \end{scope}
    \end{tikzpicture}
  \end{webinteractive}

  \webwidget[width=8cm,height=0.8cm,fallback={Select a vertex in the web deck.}]{network-selection}
\end{frame}

\begin{frame}{Canvas animation}
  \centering
  \webwidget[width=9cm,height=3.2cm,fallback={Canvas motion demo (web edition).}]{canvas-demo}
\end{frame}

\section{frames}
\begin{frame}{block samples}
    \begin{block}{block sample}
        Sample of block environment\\
        \alert{alerted text}
    \end{block} 
    \vfill
    \begin{exampleblock}{example block sample}
        Sample of example block environment
    \end{exampleblock} 
    \vfill
    \begin{alertblock}{alert block sample}
        Sample of alert block environment
    \end{alertblock} 
\end{frame}

\begin{frame}{math samples}
    \[
       \frac{1}{2\pi} \int_0^\infty f(x)e^{i\omega x} dx, \quad
       \sum_{i=1}^\infty \frac{1}{i^2} = \frac{\pi^2}{6}
    \]

    \begin{itemize}
        \item text and math \\ 
         $\alpha$-approximation algorithm, submodular function $f$, $\sum_{i=1}^n i$
        \item Blackboard bold $\mathbb{Q}$, $\mathbb{R}$, $\mathbb{C}$
        \item Bold $\mathbf{a}, \mathbf{b}, \mathbf{c}$
        \item Italic bold $\boldsymbol{a}, \boldsymbol{b}, \boldsymbol{c}$
        \item Calligraphic $\mathcal{A}$, $\mathcal{B}$, $\mathcal{C}$
        \item Operators $\max A$, $\min B$, $\sup X$, $\inf Y$
    \end{itemize}
\end{frame}

\begin{frame}{Algorithms}
    \begin{block}{Matrix Multiplication}
        \begin{algorithmic}[1]
            \STATE $C = O$
            \FOR{$i = 1, \dots, m$}
            \FOR{$j = 1, \dots, n$}
            \FOR{$k = 1, \dots, r$}
            \STATE $C[i,j] = C[i,j] + A[i, k] \cdot B[k, j]$
            \ENDFOR
            \ENDFOR
            \ENDFOR
            \RETURN $C$
        \end{algorithmic}
    \end{block}
\end{frame}
\end{document}
